\documentclass[12pt]{article}

\usepackage[margin=1in]{geometry}
\usepackage{amsmath,amssymb}
\usepackage{tikz}
\usepackage{tikz-3dplot}

\title{APM1220 FA 2}
\author{Roland Dela Rosa}
\date{Aug 2026}

\begin{document}
\maketitle

\section*{Problem 3.1}

Given

$$
X=
\begin{bmatrix}
9&1\\
5&3\\
1&2
\end{bmatrix}.
$$

\subsection*{(a)}

The three points are

$$
(9,1),\quad (5,3),\quad (1,2).
$$

The sample mean is

$$
\bar{x}
=
\begin{bmatrix}
\bar{x}_1\\
\bar{x}_2
\end{bmatrix}
=
\begin{bmatrix}
(9+5+1)/3\\
(1+3+2)/3
\end{bmatrix}
=
\begin{bmatrix}
5\\
2
\end{bmatrix}.
$$

So the sample mean is

$$
\boxed{\bar{x}=(5,2)}.
$$

\begin{center}
\begin{tikzpicture}[scale=0.65]
\draw[->] (0,0) -- (10,0) node[right] {$x_1$};
\draw[->] (0,0) -- (0,5) node[above] {$x_2$};

\foreach \x in {1,2,...,9}
\draw (\x,0.08)--(\x,-0.08) node[below] {\small \x};

\foreach \y in {1,2,3,4}
\draw (0.08,\y)--(-0.08,\y) node[left] {\small \y};

\fill (9,1) circle (3pt) node[above left] {$(9,1)$};
\fill (5,3) circle (3pt) node[above] {$(5,3)$};
\fill (1,2) circle (3pt) node[above right] {$(1,2)$};
\fill (5,2) circle (3pt) node[below right] {$\bar{x}$};
\end{tikzpicture}
\end{center}

\subsection*{(b)}

The two columns of $X$ are

$$
y_1=
\begin{bmatrix}
9\\5\\1
\end{bmatrix},
\qquad
y_2=
\begin{bmatrix}
1\\3\\2
\end{bmatrix}.
$$

Since

$$
\bar{x}_1=5,\qquad \bar{x}_2=2,
$$

we have

$$
\bar{x}_1\mathbf{1}
=
\begin{bmatrix}
5\\5\\5
\end{bmatrix},
\qquad
\bar{x}_2\mathbf{1}
=
\begin{bmatrix}
2\\2\\2
\end{bmatrix}.
$$

Thus,

$$
y_1-\bar{x}_1\mathbf{1}
=
\begin{bmatrix}
9\\5\\1
\end{bmatrix}
-
\begin{bmatrix}
5\\5\\5
\end{bmatrix}
=
\boxed{
\begin{bmatrix}
4\\0\\-4
\end{bmatrix}}
$$

and

$$
y_2-\bar{x}_2\mathbf{1}
=
\begin{bmatrix}
1\\3\\2
\end{bmatrix}
-
\begin{bmatrix}
2\\2\\2
\end{bmatrix}
=
\boxed{
\begin{bmatrix}
-1\\1\\0
\end{bmatrix}}.
$$

\begin{center}
\tdplotsetmaincoords{70}{120}
\begin{tikzpicture}[tdplot_main_coords,scale=0.65]
\draw[->] (0,0,0) -- (6,0,0) node[anchor=north east] {$x_1$};
\draw[->] (0,0,0) -- (0,6,0) node[anchor=north west] {$x_2$};
\draw[->] (0,0,0) -- (0,0,6) node[above] {$x_3$};

\draw[->] (0,0,0) -- (4,0,-4)
node[pos=.65,above] {$y_1-\bar{x}_1\mathbf{1}$};

\draw[->] (0,0,0) -- (-1,1,0)
node[pos=.6,above] {$y_2-\bar{x}_2\mathbf{1}$};

\fill (4,0,-4) circle (2.5pt);
\fill (-1,1,0) circle (2.5pt);
\end{tikzpicture}
\end{center}

\subsection*{(c)}

Let

$$
d_1=
\begin{bmatrix}
4\\0\\-4
\end{bmatrix},
\qquad
d_2=
\begin{bmatrix}
-1\\1\\0
\end{bmatrix}.
$$

The lengths are

$$
\begin{aligned}
\|d_1\|
&=\sqrt{4^2+0^2+(-4)^2}\\
&=\sqrt{32}
=4\sqrt{2},
\end{aligned}
$$

so

$$
\boxed{\|d_1\|=4\sqrt{2}}.
$$

Also,

$$
\begin{aligned}
\|d_2\|
&=\sqrt{(-1)^2+1^2+0^2}\\
&=\sqrt{2}.
\end{aligned}
$$

Therefore

$$
\boxed{\|d_2\|=\sqrt{2}}.
$$

The dot product is

$$
d_1'd_2
=
4(-1)+0(1)+(-4)(0)
=-4.
$$

Using

$$
\cos\theta
=
\frac{d_1'd_2}{\|d_1\|\|d_2\|},
$$

we get

$$
\cos\theta
=
\frac{-4}{(4\sqrt{2})(\sqrt{2})}
=-\frac12.
$$

Thus

$$
\boxed{\cos\theta=-\frac12}
$$

and

$$
\boxed{\theta=120^\circ}.
$$

The mean-corrected matrix is

$$
D=X-\mathbf{1}\bar{x}'
=
\begin{bmatrix}
4&-1\\
0&1\\
-4&0
\end{bmatrix}.
$$

Using

$$
S_n=\frac{1}{n}D'D,
$$

we get

$$
S_n
=
\frac13
\begin{bmatrix}
32&-4\\
-4&2
\end{bmatrix}
=
\begin{bmatrix}
32/3&-4/3\\
-4/3&2/3
\end{bmatrix}.
$$

The correlation matrix is

$$
\boxed{
R=
\begin{bmatrix}
1&-1/2\\
-1/2&1
\end{bmatrix}}.
$$

The lengths of the deviation vectors and the cosine of their angle are related to the diagonal and off-diagonal entries of $S_n$.

\newpage

\section*{Problem 3.6}

Given

$$
X=
\begin{bmatrix}
-1&3&-2\\
2&4&2\\
5&2&3
\end{bmatrix}.
$$

The means are

$$
\bar{x}_1=\frac{-1+2+5}{3}=2,
$$

$$
\bar{x}_2=\frac{3+4+2}{3}=3,
$$

$$
\bar{x}_3=\frac{-2+2+3}{3}=1.
$$

Therefore,

$$
\bar{x}'=
\begin{bmatrix}
2&3&1
\end{bmatrix}.
$$

\subsection*{(a)}

The deviation matrix is

$$
\begin{aligned}
D
&=X-\mathbf{1}\bar{x}'\\
&=
\begin{bmatrix}
-1&3&-2\\
2&4&2\\
5&2&3
\end{bmatrix}
-
\begin{bmatrix}
2&3&1\\
2&3&1\\
2&3&1
\end{bmatrix}\\
&=
\boxed{
\begin{bmatrix}
-3&0&-3\\
0&1&1\\
3&-1&2
\end{bmatrix}}.
\end{aligned}
$$

The rows add up to zero, so they are linearly dependent. Therefore,

$$
\operatorname{rank}(D)<3.
$$

The first two rows are independent, so

$$
\boxed{\operatorname{rank}(D)=2}.
$$

Hence $D$ is not full rank.

\subsection*{(b)}

Since $n=3$,

$$
S=\frac{1}{n-1}D'D
=\frac12D'D.
$$

First,

$$
D'D=
\begin{bmatrix}
18&-3&15\\
-3&2&-1\\
15&-1&14
\end{bmatrix}.
$$

Therefore,

$$
\boxed{
S=
\begin{bmatrix}
9&-\frac32&\frac{15}{2}\\
-\frac32&1&-\frac12\\
\frac{15}{2}&-\frac12&7
\end{bmatrix}}.
$$

Since $D$ does not have full rank, $S$ is singular. Thus,

$$
\boxed{|S|=0}.
$$

Geometrically, this means the data do not form a three-dimensional volume. They lie in a lower-dimensional space.

\subsection*{(c)}

The total sample variance is

$$
\operatorname{tr}(S)
=
9+1+7
=
\boxed{17}.
$$

\newpage

\section*{Problem 3.9}

Given

$$
X=
\begin{bmatrix}
12&17&29\\
18&20&38\\
14&16&30\\
20&18&38\\
16&19&35
\end{bmatrix}.
$$

The means are

$$
\bar{x}_1
=
\frac{12+18+14+20+16}{5}
=16,
$$

$$
\bar{x}_2
=
\frac{17+20+16+18+19}{5}
=18,
$$

$$
\bar{x}_3
=
\frac{29+38+30+38+35}{5}
=34.
$$

Thus,

$$
\bar{x}'=
\begin{bmatrix}
16&18&34
\end{bmatrix}.
$$

\subsection*{(a)}

The mean-corrected matrix is

$$
D=X-\mathbf{1}\bar{x}'
$$

$$
\boxed{
D=
\begin{bmatrix}
-4&-1&-5\\
2&2&4\\
-2&-2&-4\\
4&0&4\\
0&1&1
\end{bmatrix}}.
$$

Let the columns be $D_1,D_2,D_3$. We can see that

$$
D_3=D_1+D_2.
$$

Therefore,

$$
D_1+D_2-D_3=0.
$$

So one vector that gives the linear dependence is

$$
\boxed{
a=
\begin{bmatrix}
1\\1\\-1
\end{bmatrix}}.
$$

Indeed,

$$
Da=0.
$$

\subsection*{(b)}

The sample covariance matrix is

$$
S=\frac{1}{n-1}D'D
=\frac14D'D.
$$

We get

$$
D'D=
\begin{bmatrix}
40&12&52\\
12&10&22\\
52&22&74
\end{bmatrix}.
$$

Therefore,

$$
\boxed{
S=
\begin{bmatrix}
10&3&13\\
3&5/2&11/2\\
13&11/2&37/2
\end{bmatrix}}.
$$

Since the columns are linearly dependent,

$$
\boxed{|S|=0}.
$$

Now,

$$
Sa=
\begin{bmatrix}
10&3&13\\
3&5/2&11/2\\
13&11/2&37/2
\end{bmatrix}
\begin{bmatrix}
1\\1\\-1
\end{bmatrix}.
$$

This gives

$$
Sa=
\begin{bmatrix}
10+3-13\\
3+\frac52-\frac{11}{2}\\
13+\frac{11}{2}-\frac{37}{2}
\end{bmatrix}
=
\begin{bmatrix}
0\\0\\0
\end{bmatrix}.
$$

Thus,

$$
\boxed{Sa=0}.
$$

So $a$ is an eigenvector corresponding to the eigenvalue

$$
\boxed{\lambda=0}.
$$

\subsection*{(c)}

From the original data,

$$
29=12+17,
$$

$$
38=18+20,
$$

$$
30=14+16,
$$

$$
38=20+18,
$$

$$
35=16+19.
$$

Therefore,

$$
\boxed{x_3=x_1+x_2}.
$$

So the linear dependence is

$$
\boxed{a_1=1,\qquad a_2=1,\qquad a_3=-1}.
$$

\newpage

\section*{Problem 3.18}

The given mean is

$$
\bar{x}=
\begin{bmatrix}
0.766\\
0.508\\
0.438\\
0.161
\end{bmatrix}
$$

and the covariance matrix is

$$
S=
\begin{bmatrix}
0.856&0.635&0.173&0.096\\
0.635&0.568&0.128&0.067\\
0.173&0.127&0.171&0.039\\
0.096&0.067&0.039&0.043
\end{bmatrix}.
$$

For the calculations, I used $S_{32}=0.128$ so that the covariance matrix is symmetric.

\subsection*{(a)}

Let

$$
T=x_1+x_2+x_3+x_4.
$$

Then

$$
a=
\begin{bmatrix}
1\\1\\1\\1
\end{bmatrix}.
$$

The mean is

$$
\begin{aligned}
\bar{T}
&=a'\bar{x}\\
&=0.766+0.508+0.438+0.161\\
&=\boxed{1.873}.
\end{aligned}
$$

The variance is

$$
\operatorname{Var}(T)=a'Sa.
$$

Therefore,

$$
\begin{aligned}
\operatorname{Var}(T)
={}&0.856+0.568+0.171+0.043\\
&+2(0.635+0.173+0.096\\
&\qquad+0.128+0.067+0.039)\\
&=\boxed{3.914}.
\end{aligned}
$$

Thus,

$$
\boxed{\bar{T}=1.873,\qquad s_T^2=3.914}.
$$

\subsection*{(b)}

Let

$$
D=x_1-x_2.
$$

Then

$$
a=
\begin{bmatrix}
1\\-1\\0\\0
\end{bmatrix}.
$$

The mean is

$$
\bar{D}
=
0.766-0.508
=
\boxed{0.258}.
$$

The variance is

$$
\begin{aligned}
\operatorname{Var}(D)
&=S_{11}+S_{22}-2S_{12}\\
&=0.856+0.568-2(0.635)\\
&=\boxed{0.154}.
\end{aligned}
$$

Therefore,

$$
\boxed{\bar{D}=0.258,\qquad s_D^2=0.154}.
$$

For the covariance with the total variable,

$$
T=x_1+x_2+x_3+x_4.
$$

Using

$$
\operatorname{Cov}(a'X,b'X)=a'Sb,
$$

we get

$$
\begin{aligned}
\operatorname{Cov}(D,T)
={}&(0.856+0.635+0.173+0.096)\\
&-(0.635+0.568+0.128+0.067)\\
&=1.760-1.398\\
&=\boxed{0.362}.
\end{aligned}
$$

Thus,

$$
\boxed{\operatorname{Cov}(D,T)=0.362}.
$$

\newpage

\section*{Problem 3.20}

There are $25$ observations. Let

$$
D=x_2-x_1.
$$

\subsection*{(a)}

First calculate the means:

$$
\bar{x}_1
=
\frac{235.5}{25}
=
\boxed{9.42},
$$

$$
\bar{x}_2
=
\frac{481.8}{25}
=
\boxed{19.272}.
$$

The sample variances are

$$
s_1^2
=
\frac{339.34}{24}
=
\boxed{14.1392},
$$

and

$$
s_2^2
=
\frac{1493.7304}{24}
=
\boxed{62.2388}.
$$

The sample covariance is

$$
s_{12}
=
\frac{323.344}{24}
=
\boxed{13.4727}.
$$

The mean of the difference is

$$
\begin{aligned}
\bar{D}
&=\bar{x}_2-\bar{x}_1\\
&=19.272-9.42\\
&=\boxed{9.852}.
\end{aligned}
$$

For the variance,

$$
s_D^2
=
s_2^2+s_1^2-2s_{12}.
$$

Thus,

$$
\begin{aligned}
s_D^2
&=62.2388+14.1392-2(13.4727)\\
&=\boxed{49.4326}.
\end{aligned}
$$

Therefore,

$$
\boxed{\bar{D}=9.852,\qquad s_D^2=49.4326}.
$$

\subsection*{(b)}

Now calculate each difference $d_j=x_{2j}-x_{1j}$:

$$
\begin{array}{rrrrr}
1.2,&2.0,&9.4,&2.0,&4.1,\\
5.2,&23.1,&5.3,&4.8,&15.4,\\
11.7,&11.5,&22.5,&14.2,&8.8,\\
7.1,&5.5,&4.8,&22.6,&6.5,\\
24.8,&7.0,&9.8,&4.4,&12.6.
\end{array}
$$

Their sum is

$$
\sum d_j=246.3.
$$

Therefore,

$$
\bar{d}
=
\frac{246.3}{25}
=
\boxed{9.852}.
$$

Also,

$$
\sum d_j^2=3612.93.
$$

Using

$$
s_d^2
=
\frac{
\sum d_j^2-\dfrac{(\sum d_j)^2}{n}
}{n-1},
$$

we get

$$
\begin{aligned}
s_d^2
&=
\frac{
3612.93-\dfrac{(246.3)^2}{25}
}{24}\\
&=\boxed{49.4326}.
\end{aligned}
$$

Therefore,

$$
\boxed{\bar{d}=9.852,\qquad s_d^2=49.4326}.
$$

The results are the same as in part (a).

\end{document}
